\section{Differential calculus}

\begin{definition}
    The \textbf{average rate of change} of a function $f$ in the domain between $a$ and $a + h$ is defined as
    $$ \frac{\Delta y}{\Delta x} = \frac{f(a+h)-f(a)}{h}$$
    and also called the differential quotient.
\end{definition}

\begin{definition}
    The \textbf{derivative} $f'$ of a function $f$ at the coordinate $a$ is defined as
    $$\lim_{h \to 0}\frac{f(a+h)-f(a)}{h} = lim_{b \to a}\frac{f(b)-f(a)}{b-a} = \frac{df}{dx}\vert_{x=a} = f'(a)$$
\end{definition}

\begin{remark}
    The derivative $f'(a)$ can geometrically interpreted as the incline/decline of the graph of $f$ at the point $(a, f(a))$.
\end{remark}

\begin{definition}
    If for a function $f$ at the cooridnate $a$ the derivative $\exists f'(a)$ then f is called \textbf{differentiable} at $a$. If $f$ is differentiable for every $a \in \mathbb{D}$ then $f$ is called differentiable.
\end{definition}

\begin{example}
    let $f(x) = |x|$ we can see that for $a=0$ there is no derivative (we have a sharp change in "direction")
\end{example}

\begin{remark}
    We can generally expect that there are no derivatives at undefined points $a \notin \mathbb{D}$.
\end{remark}

\begin{lemma}
    If $f$ is differentiable at $x_0$ then $f$ is continous at $x_0$.
\end{lemma}

\begin{proof}
    Assume $f$ differentiable then
    $$\lim_{x \to x_0}f(x) = \lim_{h \to 0} f(x_0+h) = \lim_{h \to 0} f(x_0 + h) - f(x_0) + f(x_0)$$
    $$ = f(x_0) + \lim_{h \to 0} f(x_0+h)-f(x_0) = f(x_0) + \lim_{h \to 0} \frac{f(x_0+h)-f(x_0)}{h} \cdot h$$
    $$ = f(x_0) + \lim_{h \to 0} h \cdot \lim_{h \to 0} \frac{f(x_0+h)-f(x_0)}{h}$$
    Observe that the last term exists by the assumption hence
    $$ = f(x_0) + 0 \cdot f'(x_0) = f(x_0)$$
\end{proof}

\begin{definition}
    If the limit 
    $$\lim_{h>0 \to 0}\frac{f(a+h)-f(a)}{h}$$
    existst then we call it the \textbf{right-sided derivative} (or the left-sided derivative analogue)
\end{definition}

\begin{definition}
    For a function $f$ we can iteratively define higher derivatives for $n \in \mathbb{N}_0$
    $$ f^(0) = f, f^(1) = f', f^(2) = f'', \dots f^(n+1) = (f^(n))'$$
    if $f^(n)$ exists we call it $n$-times differentiable.
\end{definition}

\begin{definition}
    if $f^(n)$ is continous we call $f$ $n$-times continous differentiable. The set of all $n$-times continous differentiable functions on $D$ is denoted as $C^n(D)$
\end{definition}

\begin{definition}
    We call $C^{\infty}(D) \coloneqq \cap_{n=0}^{\infty} C^n(D)$ the set of smooth functions.
\end{definition}

\subsection{Derivation rules}

\begin{theorem}
    For functions $f(x)$ we have
    \begin{itemize}
        \item $f(x) = a x^p \quad \implies \quad f'(x) = a p x^{p-1}$
        \item $f(x) = a^x \quad \implies \quad f'(x) = ln(a) \cdot a^x$
        \item $f(x) = sin(x) \quad \implies \quad f'(x) = cos(x)$
        \item $f(x) = cos(x) \quad \implies \quad f'(x) = -sin(x)$
        \item $f(x) = tan(x) \quad \implies \quad f'(x) = \frac{1}{cos^2(x)} = 1 + tan^2(x)$
        \item $f(x) = sinh(x) \quad \implies \quad f'(x) = cosh(x)$
        \item $f(x) = cosh(x) \quad \implies \quad f'(x) = sinh(x)$
        \item $f(x) = ln(x) f'(x) = \frac{1}{x}$
        \item $f(x) = arcsin(x) f'(x) = \frac{1}{\sqrt{1-x^2}}$
        \item $f(x) = arccos(x) f'(x) = \frac{-1}{\sqrt{1-x^2}}$
        \item $f(x) = arctan(x) f'(x) = \frac{1}{1+x^2}$
    \end{itemize}
\end{theorem}

\begin{theorem}
    \textbf{Sum Rule:} For the $n$-th derivative of a sum:
    $$(f + g)^{(n)} = f^{(n)} + g^{(n)} $$
\end{theorem}

\begin{theorem}
    \textbf{General Leibniz Rule (Product):} The $n$-th derivative of a product is given by:
    $$(f \cdot g)^{(n)} = \sum_{k=0}^{n} \binom{n}{k} f^{(n-k)} g^{(k)} $$
\end{theorem}

\begin{theorem}
    \textbf{Chain Rule:} For the composition of functions:
    $$(f \circ g)'(x) = f'(g(x)) \cdot g'(x) $$
\end{theorem}

\begin{theorem}
    \textbf{Quotient Rule:} For $g(x) \neq 0$:
    $$\left(\frac{f}{g}\right)' = \frac{f'g - fg'}{g^2} $$
\end{theorem}

\begin{theorem}
    \textbf{Inverse Functions:} If $f$ is differentiable and injective:
    $$(f^{-1})'(x) = \frac{1}{f'(f^{-1}(x))} $$
\end{theorem}

\subsection{Local extrema}

\begin{theorem}
    
\end{theorem}
